% Shared, page-aligned Low-Mass course for the child field guide and trainer
% manual. Both preserve one action sequence and the same hard page boundaries;
% trainer-only meanings and correction detail belong in the adult layer.

\LMTitlePage

\BeginLMPage{PREPARATION}{VEST AND PRACTISE THE GESTURES}{%
  \LMTrainerNote{Vesture}{The Missal names the server's surplice at II.1.
    Cassock and surplice follow the manual model. Check hems and shoes before
    walking; the adult sacristan prepares candles and sacred vessels.}
  \LMTrainerNote{Reverence}{The manual's server genuflects even at an altar
    without the reserved Sacrament. The priest's different reverence does not
    change the server's. Rehearse carrying the book securely.}
  \LMTrainerNote{Source}{Fortescue, chapter IV and chapter X, sections 1--4;
    ACSS two-server sheet, pages 1--3 and 6. These practical details are the
    selected manual model.}}
\LMPageIntro{Put on your cassock, then your surplice. The cassock is the long
garment underneath; the surplice is the white garment over it. Check that your
shoes and hem let you walk and kneel safely.}
\LMModelChoice{Practise these movements with your trainer before the first
walk. Use the same gestures wherever this guide names them.}
\LMAction{1}{Both}{Join your free hands.}{Hold palms together before the chest,
fingers together and pointing up. If only one hand is busy, place the free hand
flat against the chest. Keep both hands on an object that needs two hands.}
\LMAction{2}{Both}{Know the two bows.}{A head bow lowers the head only. A
moderate bow lowers the head and shoulders together. Keep your hands joined
unless holding an object; rise smoothly after the bow.}
\LMAction{3}{Both}{Genuflect, then move.}{Stop facing the altar. Keep your body
upright, touch your right knee to the floor, and rise. Do not add a head bow or
Sign of the Cross. At arrival and departure, do this on the sanctuary floor.}
\LMAction{4}{Both}{Make the Sign of the Cross.}{Keep the left hand against the
chest. With the right hand touch forehead, breast, left shoulder, then right
shoulder. Join your hands again. For the Gospel's three small crosses, use the
right thumb on forehead, lips, and breast.}
\LMAction{5}{Both}{Begin in the sacristy.}{Before leaving, bow with the priest
to the sacristy cross. First leads, Second follows, and the priest follows
both. Carry nothing on this entrance; the Missal is already at the altar.}
\EndLMPage

% ---------------------------------------------------------------------------
% PAGE 4 — LM-002 and LM-003
% ---------------------------------------------------------------------------
\BeginLMPage{FOOT/ALTAR}{ARRIVE AND KNEEL}{%
  \LMTrainerNote{Arrival}{Fix the actual doorway and reverence before
    rehearsal. The diagram deliberately does not depict genuflection.}
  \LMTrainerNote{Biretta}{If the priest wears one, First receives and places
    it at the location fixed before Mass. Keep this out of a first rehearsal.}
  \LMTrainerNote{Watch}{Standing heads are clearly higher; kneeling figures
    show two knees on the floor and tucked feet.}}
\LMPageIntro{First stays on the priest's right. Second stays on his left.
The priest needs the centre.}
\LMModelKeyLine
\LMDiagramNote{Every position view uses the same sanctuary. Face the altar:
Gospel is picture-left, Epistle picture-right. Levels 1, 2 and 3 are the
lowest step, second step and predella. Dashed arrows show a route;
lettered stops are successive moments, not extra people.}
\vspace{4pt}
\LMSceneStart
\vspace{2pt}
\LMAction{1}{Both}{Go single file directly from the sacristy to the foot.}{First
leads and Second follows. At the foot, form side by side: First is
picture-right and Second is picture-left. Leave room for the priest.}
\LMAction{2}{Both}{Wait, then genuflect together.}{When the priest has arrived
and faced the altar, make the common reverence. Remain standing while he
places the chalice, opens the Missal, and returns to the foot.}
\EndLMPage
\BeginLMPage{FOOT/ALTAR}{WAIT, THEN KNEEL}{%
  \LMTrainerNote{Two moments}{The standing wait and opening kneel have different cues. Do not kneel merely because the priest has gone up to the altar.}
  \LMTrainerNote{Place}{Both remain on the sanctuary floor throughout the prayers at the foot.}}
\LMScenePriestPrepares
\LMRouteFields{The common arrival reverence is complete.}{Both; hands free.}{Standing at the foot.}{Stay in place while the priest prepares the altar and returns.}{Stand facing the altar, with the centre clear.}{The priest alone places the chalice and opens the Missal.}
\LMSceneFootKneel
\LMRouteFields{The priest has returned and begins the prayers.}{Both; hands free.}{The same floor positions.}{Lower together onto both knees.}{Kneel facing the altar; sign with \emph{In nomine Patris}.}{The priest's opening formula includes its Amen; add none.}
\EndLMPage

% ---------------------------------------------------------------------------
% PAGE 5 — LM-004 and first part of LM-005
% ---------------------------------------------------------------------------
\ifTrainerManualGuide\LMSuppressTrainerMeaningtrue\fi
\BeginLMPage{FOOT/ALTAR}{BEGIN THE PRAYERS}{%
  \LMTrainerNote{Pace}{First sets a moderate pace. Second begins and finishes
    with First; neither races the other.}
  \LMTrainerNote{Before 01}{Both sign with the priest at
    \emph{In nomine Patris}. His \emph{Amen} completes the formula; add none.}
  \LMTrainerNote{Branch}{At a Mass of the season from Passion Sunday through
    Holy Thursday, keep 01, omit Psalm 42, and continue at 06.}}
\LMPageIntro{Both acolytes are kneeling on the sanctuary floor at the foot.
Read across each numbered row.}
\LMBranch{Mass of the season, Passion Sunday through Holy Thursday}
{Say 01, omit 02--05 and the repeated 01, and continue at 06. Otherwise
continue with Psalm 42 below.}
\LMLaneHeadings

\LMDialogueFirst{01}
  {\emph{Introibo ad altare Dei.}}
  {A1}

\LMDialogueFirst{02}
  {\emph{Iudica me, Deus, et discerne causam meam de gente non sancta: ab
  homine iniquo, et doloso erue me.}}
  {A2}

\LMDialogueFirst{03}
  {\emph{Emitte lucem tuam, et veritatem tuam: ipsa me deduxerunt, et
  adduxerunt in montem sanctum tuum, et in tabernacula tua.}}
  {A3}

\LMReadingKey
\EndLMPage

% ---------------------------------------------------------------------------
% PAGE 6 — finish Psalm 42, LM-005 and LM-006
% ---------------------------------------------------------------------------
\BeginLMPage{FOOT/ALTAR}{FINISH PSALM 42}{%
  \LMTrainerNote{Meaning}{Use the short meanings to keep the verse chain from
    becoming unconnected sounds. They are not translations to recite.}
  \LMTrainerNote{Gesture}{Bow at the priest's \emph{Gloria Patri}. Sign
    together at \emph{Adiutorium nostrum}.}
  \LMTrainerNote{Omission}{When Psalm 42 is omitted, turn directly from 01 to
    06; do not use 02--05 or the repeated 01.}}
\LMPageIntro{Bow the head at \emph{Gloria Patri}. Make the Sign of the Cross
at \emph{Adiutorium nostrum}. Keep kneeling for both gestures.}
\LMBranch{If Psalm 42 is omitted}
{Do not use 02--05 or the repeated 01. Continue directly at 06.}
\LMLaneHeadings

\LMDialogueFirst{04}
  {\emph{Confitebor tibi in cithara, Deus, Deus meus: quare tristis es,
  anima mea, et quare conturbas me?}}
  {A4}

\LMDialogueFirst{05}
  {\emph{Gloria Patri, et Filio, et Spiritui Sancto.}}
  {A5}

\LMDialogueRepeat{01 AGAIN}
  {\emph{Introibo ad altare Dei.}}
  {A1}

\LMDialogueFirst{06}
  {\emph{Adiutorium nostrum in nomine Domini.}}
  {A6}
\EndLMPage
\ifTrainerManualGuide\LMSuppressTrainerMeaningfalse\fi

% ---------------------------------------------------------------------------
% PAGE 8 — LM-007
% ---------------------------------------------------------------------------
\BeginLMPage{FOOT/ALTAR}{ANSWER THE PRIEST'S CONFITEOR}{%
  \LMTrainerNote{Cue}{Practise the complete final sentence, not only
    \emph{Dominum Deum nostrum}.}
  \LMTrainerNote{Direction}{Both turn and make the taught bow toward the
    priest while speaking to him.}
  \LMTrainerNote{Number}{The singular words \emph{tui} and \emph{tuis}
    address one priest.}}
\LMPageIntro{Listen through the priest's complete Confiteor. His final
sentence is the cue.}
\LMLaneHeadings

\LMDialogueFirst{07}
  {\emph{Ideo precor beatam Mariam semper Virginem, beatum Michaelem
  Archangelum, beatum Ioannem Baptistam, sanctos Apostolos Petrum et Paulum,
  omnes Sanctos, et vos, fratres, orare pro me ad Dominum Deum nostrum.}}
  {A7}

\LMAction{1}{Both}{Turn and bow toward the priest.}{Say the whole response
together. The priest answers \emph{Amen}. Turn back toward the altar, remain
bowed, and begin your
Confiteor in Part 1 below.}
\EndLMPage

% ---------------------------------------------------------------------------
% PAGES 9--11 — LM-008, complete Confiteor phrase ladder
% ---------------------------------------------------------------------------
\ContinueLMPairedScene{FOOT/ALTAR}{YOUR CONFITEOR \MissalSeparator\ PART 1}{%
  \LMTrainerNote{One prayer}{The three parts divide one complete response for
    learning. Join them without a pause once the prayer is learned.}
  \LMTrainerNote{Turn}{At \emph{et tibi, pater}, both turn toward the priest,
    then face the altar again.}
  \LMTrainerNote{Deck}{The full Confiteor is intentionally absent from the
    fixed-size eight-up flash cards; shrinking it would recreate the tested
    clutter problem.}}
\LMPageIntro{Both remain bowed. Learn one sense-sized part at a time, then
join all three parts without stopping. Turn toward the priest at
\emph{et tibi, pater}, then face the altar again.}
\LMPhraseBlock{PART 1}{I CONFESS}
  {\AightOneLatin}
  {\AightOneSound}
  {\AightOneMeaning}
\LMAction{1}{Both}{Keep the prayer moving.}{At the colon after
\emph{pater}, continue directly to Part 2.}
\EndLMPage

\BeginLMPage{FOOT/ALTAR}{YOUR CONFITEOR \MissalSeparator\ PART 2}{%
  \LMTrainerNote{Breast}{At each \emph{mea culpa}, both strike the breast
    once with the right hand.}
  \LMTrainerNote{Endings}{Correct the Latin word endings before increasing
    speed.}
  \LMTrainerNote{Join}{Begin at Part 1, continue through this page, then turn
    immediately to Part 3.}}
\LMPageIntro{Part 2 begins with the reason for the confession and starts the
petition for prayer.}
\LMPhraseBlock{PART 2}{BECAUSE I SINNED}
  {\AightTwoLatin}
  {\AightTwoSound}
  {\AightTwoMeaning}
\LMAction{2}{Both}{Strike three times.}{Strike once at each
\emph{mea culpa}. Keep the other hand over the breast.}
\EndLMPage

\ContinueLMPairedScene{FOOT/ALTAR}{YOUR CONFITEOR \MissalSeparator\ PART 3}{%
  \LMTrainerNote{Turn}{At \emph{et te, pater}, both turn toward the priest,
    then finish still bowed.}
  \LMTrainerNote{Whole run}{Once each part is secure, start at
    \emph{Confiteor Deo} and finish at \emph{Deum nostrum} without a prompt.}
  \LMTrainerNote{Next}{The priest now says \emph{Misereatur vestri}; both
    remain bowed for the first Amen.}}
\LMPageIntro{Part 3 completes the petition. This is the end of the one full
Confiteor begun in Part 1.}
\LMPhraseBlock{PART 3}{PLEASE PRAY FOR ME}
  {\AightThreeLatin}
  {\AightThreeSound}
  {\AightThreeMeaning}
\LMAction{3}{Both}{Turn toward the priest at \emph{et te, pater}.}{Face the
altar again and finish together. Stay bowed after
\emph{Dominum Deum nostrum}. Listen for the priest's next prayer.}
\EndLMPage

% ---------------------------------------------------------------------------
% PAGES 12--13 — LM-009 and LM-010
% ---------------------------------------------------------------------------
\BeginLMPage{FOOT/ALTAR}{RECEIVE THE ABSOLUTION PRAYERS}{%
  \LMTrainerNote{Posture}{Both remain bowed through the first Amen. At
    \emph{Indulgentiam}, rise upright and sign with the priest.}
  \LMTrainerNote{Amen}{This is the first taught Amen. Use the same complete
    word at later audible prayer endings.}
  \LMTrainerNote{No extra}{Do not carry this Amen backward to the opening
    Sign of the Cross.}}
\LMPageIntro{Stay bowed through the first Amen. At \emph{Indulgentiam}, become
upright on your knees and make the Sign of the Cross. Answer Amen after each
prayer; then bow the head at \emph{Deus, tu conversus}.}
\LMLaneHeadings

\LMDialogueFirst{08}
  {\emph{Misereatur vestri omnipotens Deus, et, dimissis peccatis vestris,
  perducat vos ad vitam aeternam.}}
  {A9}

\LMDialogueRepeat{08 AGAIN}
  {\emph{Indulgentiam, absolutionem et remissionem peccatorum nostrorum
  tribuat nobis omnipotens et misericors Dominus.}}
  {A9}

\LMDialogueFirst{09}
  {\emph{Deus, tu conversus vivificabis nos.}}
  {A10}
\EndLMPage

\BeginLMSingleScenePage{FOOT/ALTAR}{FINISH AT THE FOOT}{%
  \LMTrainerNote{Bow}{Both make the taught simple bow through the three short
    versicles, then raise their heads at \emph{Oremus}.}
  \LMTrainerNote{Alb}{If the local route assigns it, both assist the priest's
    alb as he ascends.}
  \LMTrainerNote{Transition}{After \emph{Oremus}, stand and move directly to
    the normal places shown on the next page.}}
\LMPageIntro{Keep the head bowed through the short answers and greeting.
Raise it at \emph{Oremus}; stand when the priest begins to ascend.}
\LMLaneHeadings

\LMDialogueFirst{10}
  {\emph{Ostende nobis, Domine, misericordiam tuam.}}
  {A11}

\LMDialogueFirst{11}
  {\emph{Domine, exaudi orationem meam.}}
  {A12}

\LMDialogueFirst{12}
  {\emph{Dominus vobiscum.}}
  {A13}
\EndLMPage

% ---------------------------------------------------------------------------
% PAGE 14 — LM-010 and LM-011
% ---------------------------------------------------------------------------
\BeginLMPage{READINGS}{MOVE TO THE NORMAL PLACES}{%
  \LMTrainerNote{Cue}{Both raise their heads at \emph{Oremus}; rise only after the priest starts ascending.}
  \LMTrainerNote{Levels}{The normal places are on Step 1, two levels below the predella. Map this model to the actual sanctuary.}
  \LMTrainerNote{Introit}{Both sign with the priest at the Introit opening and bow the head at \emph{Gloria Patri}.}}
\LMSceneRiseFromFoot
\LMRouteFields{The priest begins to ascend.}{Both; hands free.}{Kneeling at the foot on the floor.}{Rise, turn inward, and take the short outside routes.}{Reach the normal places facing the altar.}{Keep the priest's centre lane free.}
\LMSceneNormalPlaces
\LMAction{1}{Both}{Kneel together on Step 1.}{Second is Gospel, picture-left; First is Epistle, picture-right. Remain for the Introit, Kyrie, Gloria when appointed, Collect, and Epistle. Sign at the Introit opening and bow at its \emph{Gloria Patri}.}
\EndLMPage

% ---------------------------------------------------------------------------
% PAGE 15 — LM-012 and LM-013
% ---------------------------------------------------------------------------
\BeginLMPage{READINGS}{SAY THE NINE KYRIES}{%
  \LMTrainerNote{Pattern}{Low Mass uses three groups of three. Both acolytes
    speak every server position together.}
  \LMTrainerNote{Cards}{Flash cards 13 and 14 teach the two priest-to-server
    pair types; only this page shows who begins each group.}
  \LMTrainerNote{Gloria}{If the Gloria follows, remain kneeling and make only
    the gestures taught locally; there is no short server response.}}
\LMPageIntro{The priest and both acolytes alternate. Do not turn the two
\emph{Christe} server places into priest cues.}

\vspace{4pt}
{\fontsize{8}{9.5}\selectfont\bfseries GROUP 1 \textperiodcentered\ KYRIE\par}
\vspace{3pt}
{\fontsize{13}{16}\selectfont
Priest: \emph{Kyrie, eleison.}\quad
\textbf{Both: \emph{Kyrie, eleison.}}\quad
Priest: \emph{Kyrie, eleison.}\par}
\vspace{8pt}\hrule height 0.5pt
\vspace{8pt}
{\fontsize{8}{9.5}\selectfont\bfseries GROUP 2 \textperiodcentered\ CHRISTE\par}
\vspace{3pt}
{\fontsize{13}{16}\selectfont
\textbf{Both: \emph{Christe, eleison.}}\quad
Priest: \emph{Christe, eleison.}\quad
\textbf{Both: \emph{Christe, eleison.}}\par}
\vspace{8pt}\hrule height 0.5pt
\vspace{8pt}
{\fontsize{8}{9.5}\selectfont\bfseries GROUP 3 \textperiodcentered\ KYRIE\par}
\vspace{3pt}
{\fontsize{13}{16}\selectfont
Priest: \emph{Kyrie, eleison.}\quad
\textbf{Both: \emph{Kyrie, eleison.}}\quad
Priest: \emph{Kyrie, eleison.}\par}

\vspace{12pt}
{\fontsize{7.5}{9}\selectfont\bfseries SAY IT\par}
\vspace{2pt}
{\fontsize{12}{14.5}\selectfont
\ServerSoundData{A14}\par}
\vspace{7pt}
{\fontsize{7.5}{9}\selectfont\bfseries MEANING\par}
\vspace{2pt}
{\fontsize{12}{14.5}\selectfont
\ServerMeaningData{A14}\par}
\EndLMPage

% ---------------------------------------------------------------------------
% PAGE 16 — LM-014 and LM-015
% ---------------------------------------------------------------------------
\BeginLMPage{READINGS}{COLLECT AND EPISTLE}{%
  \LMTrainerNote{Collects}{Answer Amen at the audible ending of each appointed
    Collect.}
  \LMTrainerNote{Lessons}{On days with several lessons, answer only where the
    proper sheet appoints it; move the Missal only after the final lesson.}
  \LMTrainerNote{Example}{The complete Epistle cue is a verified exemplar
    from the Fifth Sunday after Pentecost, not the ending for every day.}}
\LMPageIntro{Remain kneeling. The first two responses are repeats. Card 15
uses one complete Epistle ending for practice.}
\LMLaneHeadings

\LMDialogueRepeat{12 AGAIN}
  {\emph{Dominus vobiscum.}}
  {A13}

\LMDialogueRepeat{08 AGAIN}
  {\emph{Per omnia saecula saeculorum.}}
  {A9}

\LMDialogueFirst{15}
  {\emph{Dominum autem Christum sanctificate in cordibus vestris.}}
  {A15}
\EndLMPage

% ---------------------------------------------------------------------------
% PAGE 17 — LM-016
% ---------------------------------------------------------------------------
\BeginLMPage{READINGS}{TAKE THE MISSAL AROUND THE FOOT}{%
  \LMTrainerNote{Book}{First waits until the priest clears it. Grip the stand securely at both sides; never pull its front edge or reach around the priest.}
  \LMTrainerNote{Second}{Second stands at his normal Gospel place. He does not climb to First's level. Step outward only enough to leave the passing lane clear.}
  \LMTrainerNote{Model}{The First-acolyte book assignment follows the selected O'Connell model. Keep the assigned roles throughout.}}
\LMSceneMissalMove
\LMSceneMissalCentre
\LMRouteFields{After the final \emph{Deo gratias}, the priest clears the book.}{First carries the Missal and stand; Second waits.}{First at Epistle; Second at his Gospel normal place.}{First takes the book down Epistle, stops at the centre for the taught reverence, then goes up Gospel.}{Place the book angled toward the priest; First steps down to Step 2 beside it.}{Second stands at his normal place and clears First's route without climbing.}
\EndLMPage

% ---------------------------------------------------------------------------
% PAGE 18 — LM-017
% ---------------------------------------------------------------------------
\BeginLMPage{READINGS}{STAY FOR THE GOSPEL OPENING}{%
  \LMTrainerNote{Two levels}{First waits at the left/Gospel side on Step 2. Second remains standing at his normal Gospel place on the floor, below the steps.}
  \LMTrainerNote{Timing}{First does not leave during the greeting or title. He waits until the priest has begun the reading and completed any immediate bow.}
  \LMTrainerNote{Crosses}{Both use the right thumb on forehead, lips, and breast.}}
\LMSceneGospel
\LMPageIntro{The priest faces left to read. Both stand and answer. At the title, make the three small crosses. First stays by the book until the reading begins. Second stays at his normal place below the steps.}
\LMLaneHeadings
\LMSuppressTrainerMeaningtrue
\LMDialogueRepeat{12 AGAIN}{\emph{Dominus vobiscum.}}{A13}
\LMDialogueFirst{16}{\emph{Sequentia sancti Evangelii secundum Matthaeum.}}{A16}
\LMSuppressTrainerMeaningfalse
\EndLMPage

\BeginLMPage{READINGS}{RETURN AFTER THE READING BEGINS}{%
  \LMTrainerNote{Cue}{The title response and crosses are complete. First waits until the text has begun, normally through \emph{In illo tempore} and the immediate reverence.}
  \LMTrainerNote{Ending}{Both answer \emph{Laus tibi, Christe} when the reading ends; do not wait for the priest to kiss the book.}}
\LMSceneGospelReturn
\LMRouteFields{The priest has begun the Gospel text and completed any immediate reverence.}{First, empty-handed.}{Gospel Step 2, beside the book.}{Descend below the steps, stop at the centre for the taught reverence, and return to the Epistle normal place.}{Stand on the floor below the steps, turned toward the priest; Second stays at his normal Gospel place.}{Second leaves the lane clear. Both join any appointed Gospel genuflection.}
\LMLaneHeadings
\LMSuppressTrainerMeaningtrue
\LMDialogueFirst{17}{\emph{Relinque ibi munus tuum ante altare, et vade prius reconciliari fratri tuo: et tunc veniens offeres munus tuum.}}{A17}
\LMSuppressTrainerMeaningfalse
\EndLMPage

% ---------------------------------------------------------------------------
% PAGE 19 — LM-018 and LM-019
% ---------------------------------------------------------------------------
\BeginLMPage{READINGS}{AFTER THE GOSPEL}{%
  \LMTrainerNote{Sermon}{Fix the actual seating place and return signal before
    rehearsal. Keep First and Second together.}
  \LMTrainerNote{Creed}{This core model keeps both kneeling. Teach the
    appointed reverence at \emph{Et incarnatus est} live; it is not pictured.}
  \LMTrainerNote{No words}{Neither branch adds a short server response.}}
\LMPageIntro{After \emph{Laus tibi, Christe}, follow the branch for this
Mass. Then prepare for the Offertory.}
\LMCompactTrainerAction{1}{Both}{If there is a sermon: follow the trained route.}{Go
together to the place fixed by the trainer and sit. At the signal, return
together and kneel at the normal places.}
\LMCompactTrainerAction{2}{Both}{If the Creed is appointed: kneel.}{Keep
hands joined. At \emph{Et incarnatus est} through \emph{et homo factus est},
bow the head and shoulders while remaining on both knees. Add no response.}

\LMBranch{After either branch}
{Whether or not there was a sermon, and whether or not the Creed is appointed,
return to the normal places, kneel together, and watch for the priest to begin
the Offertory.}
\EndLMPage

% ---------------------------------------------------------------------------
% PAGE 20 — LM-020 and LM-021
% ---------------------------------------------------------------------------
\BeginLMPage{OFFERTORY}{RECEIVE AND KEEP THE VEIL SAFE}{%
  \LMTrainerNote{Cue}{There is no response after \emph{Oremus}. When the priest removes the chalice veil, both rise; First then approaches the rehearsed receiving point.}
  \LMTrainerNote{Veil}{First leaves the folded veil at its safe Epistle pickup point for the later exchange. Neither server takes the chalice.}
  \LMTrainerNote{Objects}{First takes wine and Second takes water. Open the cruets at the credence before approaching.}}
\LMPageIntro{Answer the greeting and bow the head at \emph{Oremus}. Remain kneeling until the priest removes the chalice veil.}
\LMLaneHeadings
\LMDialogueRepeat{12 AGAIN}{\emph{Dominus vobiscum.}}{A13}
\LMSceneVeilReceipt
\LMRouteFields{The priest removes the chalice veil.}{First receives the folded veil; Both rise.}{The normal kneeling places.}{First approaches the priest's rehearsed Epistle receiving point and takes the veil.}{Place it folded at the safe Epistle pickup point for the later exchange.}{Second stands at his normal Gospel place until First is ready; neither takes the chalice.}
\EndLMPage

\BeginLMPage{OFFERTORY}{GO FOR THE CRUETS}{%
  \LMTrainerNote{Start}{The folded veil is secure. First returns from its pickup point; Second joins him at the foot.}
  \LMTrainerNote{Objects}{First takes wine and Second water. Open the cruets at the credence and keep the handles free for the priest.}}
\LMSceneCruetPickup
\LMRouteFields{First has secured the folded veil.}{Both; First takes wine, Second water.}{First at the Epistle veil point; Second at his standing Gospel place.}{Meet at the foot, make the taught reverence, and go to the credence.}{Take the open cruets, then stand side by side on Step 2 facing the priest.}{The folded veil stays at its safe pickup point; the priest remains on the predella.}
\EndLMPage

\BeginLMPage{OFFERTORY}{PRESENT THE WINE HANDLE FIRST}{%
  \LMTrainerNote{Kiss}{Carmody pp. 70--72 places the light kiss at the base of the handle. Never kiss the spout or the priest's hand.}
  \LMTrainerNote{Grip}{Support the cruet securely with the handle turned to the priest. Keep the other cruet upright and still.}
  \LMTrainerNote{Return}{Receive the wine back without another kiss. Second still waits for the water blessing.}}
\LMSceneWinePresentation
\LMCruetSequenceDetail
\LMRouteFields{Both have reached the service station and bowed.}{First: wine; Second: water.}{Side by side on Step 2, one level below the priest.}{First lightly kisses the base of the wine-cruet handle, then turns the handle toward the priest and offers it within his reach.}{Receive the same cruet back and hold it upright.}{Second holds water and waits; the priest takes the wine by its handle.}
\EndLMPage

\BeginLMPage{OFFERTORY}{WAIT, THEN PRESENT THE WATER}{%
  \LMTrainerNote{Order}{Wait for the priest's blessing of the water. Then Second kisses the base of its handle and presents it.}
  \LMTrainerNote{Return}{The selected Carmody sequence prescribes no second kiss on either returned cruet.}
  \LMTrainerNote{Next}{After both cruets are returned, bow together and go to the credence for the Lavabo articles.}}
\LMSceneWaterPresentation
\LMCruetPresentationDetail
\LMRouteFields{The priest has blessed the water.}{Second: water; First holds the returned wine.}{The same side-by-side service station.}{Second lightly kisses the base of the handle, turns the handle toward the priest, and presents the water.}{Receive it back without another kiss; both bow, turn inward, and return to the credence.}{First holds the wine upright and leaves Second's handoff clear.}
\EndLMPage

% ---------------------------------------------------------------------------
% PAGE 21 — LM-022 and LM-023
% ---------------------------------------------------------------------------
\BeginLMPage{OFFERTORY}{WASH AND DRY THE PRIEST'S FINGERS}{%
  \LMTrainerNote{Places}{Second brings water and bowl; First brings the opened towel. Both use Step 2 outside the Epistle edge.}
  \LMTrainerNote{Water}{A small controlled stream falls over the fingers into the bowl. This is not the later washing into the chalice.}
  \LMTrainerNote{Towel}{Keep the towel clear of the water. Present it, receive it back, then bow together.}}
\LMSceneLavaboAction
\LMLavaboActionDetail
\LMRouteFields{The cruets are reset and the priest comes to the Lavabo place.}{Second: water and bowl; First: towel.}{The credence.}{Bring the articles to the outside Epistle station on Step 2 and bow.}{Second pours into the bowl; First presents the open towel and receives it back. Face the priest.}{Each keeps the other server's hands clear; bow together when finished.}
\EndLMPage

\BeginLMPage{OFFERTORY}{RESET, THEN BRING THE BELL}{%
  \LMTrainerNote{Reset}{Return all Lavabo articles. Replace the cruet stoppers loosely for the later ablutions.}
  \LMTrainerNote{Bell}{First alone takes it. Hold it level, keep it from striking, and set it silently on Step 2.}}
\LMSceneLavaboReset
\LMSceneBell
\LMRouteFields{The priest has finished drying his fingers.}{Both reset the articles; First carries the bell.}{The Lavabo service station.}{Return to the credence; then meet at the foot, genuflect, and go to the normal places.}{Both kneel facing the altar; First places the bell on Step 2 in front of him.}{Second never reaches across First for the bell.}
\EndLMPage

% ---------------------------------------------------------------------------
% PAGE 22 — LM-024
% ---------------------------------------------------------------------------
\BeginLMPage{OFFERTORY}{ANSWER ORATE, FRATRES}{%
  \LMTrainerNote{Cue}{Wait until the priest has completed
    \emph{Orate, fratres} and turned back to the altar.}
  \LMTrainerNote{Word}{The server says \emph{de manibus tuis}. The priest's
    exceptional self-answer \emph{meis} is not a server form.}
  \LMTrainerNote{Place}{Both remain kneeling at the normal places.}}
\LMPageIntro{After the Lavabo, both are kneeling at the normal places. Listen
through the priest's complete invitation before answering.}
\LMLaneHeadings

\LMDialogueFirst{18}
  {\emph{Orate, fratres: ut meum ac vestrum sacrificium acceptabile fiat
  apud Deum Patrem omnipotentem.}}
  {A18}
\EndLMPage

% ---------------------------------------------------------------------------
% PAGE 23 — LM-025
% ---------------------------------------------------------------------------
\ContinueLMPairedScene{CANON}{ANSWER THE PREFACE DIALOGUE}{%
  \LMTrainerNote{Stage}{The Preface dialogue leads directly into the Canon,
    so it is kept in this learner stage.}
  \LMTrainerNote{Bow}{Both make a simple head bow when the priest says
    \emph{Gratias agamus Domino Deo nostro}.}
  \LMTrainerNote{Voice}{All four responses are spoken by Both in this
    two-server Low-Mass model.}}
\LMPageIntro{Both kneel at the normal places. Bow the head at
\emph{Gratias agamus Domino Deo nostro}. First's bell rests on the second
step in front of him.}
\LMLaneHeadings

\LMDialogueRepeat{08 AGAIN}
  {\emph{Per omnia saecula saeculorum.}}
  {A9}

\LMDialogueRepeat{12 AGAIN}
  {\emph{Dominus vobiscum.}}
  {A13}

\LMDialogueFirst{19}
  {\emph{Sursum corda.}}
  {A19}

\LMDialogueFirst{20}
  {\emph{Gratias agamus Domino Deo nostro.}}
  {A20}
\EndLMPage

% ---------------------------------------------------------------------------
% PAGE 24 — LM-026 and LM-027
% ---------------------------------------------------------------------------
\BeginLMPage{CANON}{FOLLOW THE BELL PLAN}{%
  \LMTrainerNote{Pattern}{This teaching model uses three clear Sanctus rings,
    one warning ring, and three distinct rings at each elevation.}
  \LMTrainerNote{Local}{The Missal permits continuous elevation ringing.
    Use only the pattern selected by the priest for the actual church.}
  \LMTrainerNote{Pickup}{Lift and set down the bell silently. Let the sound
    fade before moving it.}}
\LMPageIntro{First handles the bell. Second watches First's movement cue.
Both remain kneeling until the warning ring is complete.}
\LMModelChoice{Use only the priest's selected bell pattern. This guide uses
three Sanctus rings, one warning, and three distinct rings at each elevation.}
\LMSceneBell
\vspace{4pt}
\LMDiagramNote{The bell rests on the \textbf{second step} in front of First.
Second never reaches across First for it.}

\LMAction{1}{First}{At \emph{Sanctus}: three clear rings.}{Set the bell down
quietly. Both remain kneeling at the normal places.}
\LMAction{2}{First}{Shortly before the Consecration: one warning ring.}{Use
the priest's hands extended over the offerings as the rehearsed visual cue.}
\LMAction{3}{Both}{Rise and ascend the short way.}{Kneel on the predella,
one on each side and slightly behind the priest. First carries the bell.}
\EndLMPage

% ---------------------------------------------------------------------------
% PAGE 25 — LM-028 and LM-029
% ---------------------------------------------------------------------------
\BeginLMPage{CANON}{ASSIST AT THE HOST ELEVATION}{%
  \LMTrainerNote{Nearest edge}{Second uses his right hand and First his left to lift the nearest lower edge. Do not pull sideways or touch the priest's back.}
  \LMTrainerNote{Bell}{First rings with his outside right hand. This model uses three separate rings at each elevation.}}
\LMSceneHostElevation
\LMHostElevationDetail
\LMRouteFields{First's warning ring is complete; the Host elevation follows.}{Both; First carries the bell and rings.}{The normal kneeling places.}{Rise, ascend the short way, and kneel slightly behind the priest. When he raises the Host, gently lift the nearest lower edge; First rings three distinct times.}{Release as the Host is lowered; remain kneeling facing the altar.}{Keep the bell clear of the vestment and wait for the Chalice elevation.}
\EndLMPage

\BeginLMPage{CANON}{ASSIST AT THE CHALICE ELEVATION}{%
  \LMTrainerNote{Repeat}{Keep the same places and nearest-edge hands. First rings three distinct times again.}
  \LMTrainerNote{Return cue}{Release the chasuble when the Chalice is lowered. Wait for the priest's final reverence before rising to leave.}}
\LMSceneChaliceElevation
\LMChaliceElevationDetail
\LMRouteFields{The priest raises the Chalice.}{Both lift the chasuble; First rings.}{The same kneeling places on the predella.}{Repeat the gentle lift and three separate rings.}{Release as the Chalice is lowered; stay facing the altar.}{Both remain still until the priest completes his final reverence.}
\EndLMPage
\BeginLMPage{CANON}{RETURN AFTER THE ELEVATIONS}{%
  \LMTrainerNote{Cue}{Wait until the priest has completed his final reverence.}
  \LMTrainerNote{Bell}{First carries the bell quietly. The descent and centre meeting are rehearsed with empty hands first.}}
\LMPageIntro{Both start kneeling on the predella. The priest remains at the altar.}
\LMSceneElevationReturn
\LMAction{3}{Both}{Return to the normal places.}{Rise, face inward, descend
the short way, genuflect at the foot, and kneel on the first
step. First carries the bell without ringing it and sets it silently on Step 2.}
\ASGMovementFields{The priest has finished his final reverence.}{Both; First carries the bell.}{Descend the short way and meet at the centre stop A.}{After the centre genuflection, First kneels Epistle and Second Gospel on Step 1.}{The priest stays at the altar; neither server adds a bell signal.}
\EndLMPage

% ---------------------------------------------------------------------------
% PAGE 26 — LM-030 and LM-031
% ---------------------------------------------------------------------------
\BeginLMPage{COMMUNION}{ANSWER AFTER THE CANON}{%
  \LMTrainerNote{Quiet Amen}{After \emph{Sed libera nos a malo}, the next
    quiet Amen belongs to the priest. Do not add one.}
  \LMTrainerNote{Agnus Dei}{There is no short server response. Both remain
    kneeling and make the gestures taught for its three petitions.}
  \LMTrainerNote{Pax}{An actual pax-instrument branch is not part of this core
    route; train it separately if appointed.}}
\LMPageIntro{Both are kneeling at the normal places. The familiar Amen and
greeting responses return.}
\LMLaneHeadings

\LMDialogueRepeat{08 AGAIN}
  {\emph{Per omnia saecula saeculorum.}}
  {A9}

\LMDialogueFirst{21}
  {\emph{Et ne nos inducas in tentationem.}}
  {A21}

\LMDialogueRepeat{08 AGAIN}
  {\emph{Per omnia saecula saeculorum.}}
  {A9}

\LMDialogueRepeat{12 AGAIN}
  {\emph{Pax Domini sit semper vobiscum.}}
  {A13}
\EndLMPage

% ---------------------------------------------------------------------------
% PAGE 27 — LM-032 through LM-035
% ---------------------------------------------------------------------------
\BeginLMPage{COMMUNION}{FOLLOW THE COMMUNION BRANCH}{%
  \LMTrainerNote{1962 correction}{Do not insert a second Confiteor,
    Misereatur, or Indulgentiam before the people's Communion.}
  \LMTrainerNote{Warning}{If the faithful will receive, this model rings once
    after each of the priest's three \emph{Domine non sum dignus}.}
  \LMTrainerNote{Plate bearer}{In this model, First alone handles the
    Communion plate. First sets down the bell before taking it; Second never
    receives it.}
  \LMTrainerNote{After Communion}{First gives the plate to the priest for
    purification. It stays at the altar until the final wine and water.}}
\LMPageIntro{There are no server words at the priest's
\emph{Domine non sum dignus}, at \emph{Ecce Agnus Dei}, or at the formula
used to give Communion.}
\LMModelChoice{If Communion is given, this guide assigns First the plate and
three warning rings. The priest or trainer must mark the side and lane.}

\LMAction{1}{Both}{At the Agnus Dei: remain kneeling.}{Make the taught bow
and the three breast strikes at the three petitions. Add no response.}
\LMAction{2}{First}{If the faithful will receive: ring three times.}{Ring once
after each of the priest's three \emph{Domine non sum dignus}. If nobody will
receive, omit this Communion warning. Set the bell down silently.}
\LMAction{3}{First}{If the faithful will receive: rise and take the Communion
plate.}{Take it from the credence. Join the priest on the side marked by the
trainer. Keep the
plate level and under each communicant's chin. If nobody will receive, leave
the plate at the credence and remain in place.}
\LMScenePlatePickup
\EndLMPage

\BeginLMPage{COMMUNION}{FOLLOW WITH THE COMMUNION PLATE}{%
  \LMTrainerNote{Side}{The model assigns First the priest's right side. Map the real rail and moving lane before using objects.}
  \LMTrainerNote{Level}{Keep the plate horizontal under the chin; move after the Host has been given, not ahead of the priest.}
  \LMTrainerNote{Fragments}{Only the priest handles any fragments. First returns the plate after the priest has cleared it and gives the signal.}}
\LMPageIntro{This page applies when Communion is given to the faithful. First alone carries the plate; Second remains at the normal Gospel place.}
\LMScenePlate
\vspace{3pt}
\LMDiagramNote{\textbf{MODEL SHOWN: priest's right, picture-left.}
Use the side and lane marked at the actual church. First follows with the
plate; Second remains clear. After the last communicant, First gives the plate to the priest. The
priest clears it at the altar; First retrieves it after the final wine and
water.}
\LMRouteFields{The priest begins distributing Communion.}{First; Communion plate.}{Beside the priest, at the rehearsed side.}{Follow him along the Communion route, keeping the plate level under each chin.}{After the last communicant, give the plate to the priest at the taught altar point for purification.}{Second remains clear at Gospel; First does not take the plate back until it has been cleared.}

\LMBranch{If nobody receives}
{Leave the plate at the credence. When the priest offers the chalice for
purification, continue directly to the ablutions.}
\EndLMPage

% ---------------------------------------------------------------------------
% PAGE 28 — LM-036
% ---------------------------------------------------------------------------
\BeginLMPage{COMMUNION}{FIRST BRINGS BOTH CRUETS}{%
  \LMTrainerNote{Cue}{With Communion: move when the tabernacle is closed. Without Communion: move when the priest offers the chalice.}
  \LMTrainerNote{One server}{First serves all three pours. Second stays at the Gospel-side foot and prepares for the later book move.}
  \LMTrainerNote{Separate acts}{Wine first purifies the chalice. Only afterward are wine and water poured over the fingers.}}
\LMSceneAblutionPickup
\LMRouteFields{The tabernacle is closed, or the priest offers the chalice when there was no distribution.}{First carries both cruets.}{First at the Communion return or normal place.}{Go by the rehearsed centre route to the credence, then approach Epistle.}{Stand on Step 2 at Epistle, picture-right, holding both cruets upright.}{Second rises when First leaves for the cruets, then stands clear at the Gospel-side foot; the priest stays on the predella.}
\EndLMPage

\BeginLMPage{COMMUNION}{POUR THE FIRST WINE PURIFICATION}{%
  \LMTrainerNote{Chalice}{For this first pour the priest extends and holds the chalice. It is separate from the later finger washing with the chalice on the altar.}
  \LMTrainerNote{Signal}{First ascends to the predella, pours the wine, and stops when the priest raises the chalice. Return to Step 2 and wait.}}
\LMSceneWinePurification
\LMWinePurificationDetail
\LMRouteFields{The priest extends the chalice for wine.}{First alone; wine ready, water held upright.}{The Epistle service station.}{Ascend to the predella and approach the priest at the centre; pour wine into the priest-held chalice and stop at his signal.}{Return to Step 2 at Epistle, picture-right, facing him.}{Wait while the priest drinks this purification and comes for the finger washing; Second stays clear at Gospel.}
\EndLMPage

\BeginLMPage{COMMUNION}{WASH THE FINGERS WITH WINE}{%
  \LMTrainerNote{Altar contact}{The chalice now rests on the altar. The priest holds his joined thumbs and index fingers over its bowl.}
  \LMTrainerNote{Levels}{The priest is on the predella; First is at Epistle, picture-right, one step below. First alone pours.}}
\LMSceneFingerWine
\LMFingerWineDetail
\LMRouteFields{The priest has drunk the first purification and presents his fingers over the chalice.}{First alone; wine.}{Step 2 at Epistle, picture-right.}{Remain there and pour a little wine over the joined thumbs and index fingers.}{Stop at the priest's signal; the liquid falls into the chalice resting on the altar.}{Keep the water ready; Second remains clear at Gospel.}
\EndLMPage

\BeginLMPage{COMMUNION}{FOLLOW WITH WATER, THEN RESET}{%
  \LMTrainerNote{Sequence}{Wine over the fingers is followed by water over the same fingers. The priest stays on the predella, the chalice on the altar.}
  \LMTrainerNote{Plate}{After returning both cruets, First retrieves any Communion plate only when the priest has cleared it and signals.}}
\LMSceneFingerWater
\LMFingerWaterDetail
\LMRouteFields{The wine has finished.}{First alone; water.}{The same Step 2 station.}{Pour water over the same fingers into the chalice; stop at the signal.}{Return both cruets to the credence. Return any cleared Communion plate level after the priest's signal.}{Second stays clear and ready to move the Missal; First leaves the chalice to the priest.}
\EndLMPage

% ---------------------------------------------------------------------------
% PAGE 29 — LM-037
% ---------------------------------------------------------------------------
\BeginLMPage{COMMUNION}{EXCHANGE THE VEIL AND MISSAL}{%
  \LMTrainerNote{Order}{Begin only after First has reset both cruets and any
    used plate. First takes the folded chalice veil from Epistle to Gospel;
    Second takes the Missal and stand from Gospel to Epistle.}
  \LMTrainerNote{Book state}{The Missal ends straight on the Epistle side,
    as at the Introit.}
  \LMTrainerNote{Chalice}{First carries only the folded ornamental chalice
    veil, never the chalice. The priest alone reassembles and covers the
    chalice. The acolytes' return cue is completion of their object placements
    and centre reverence, not an invented universal ``chalice covered'' cue.}}
\LMPageIntro{First and Second change sides with separate objects: First
carries the folded chalice veil from Epistle to Gospel while Second carries
the Missal and stand from Gospel to Epistle. After both objects are placed,
they genuflect at the centre and switch back to their normal
positions. Fix the exact lanes at the actual altar.}
\LMSceneBookReturn
\vspace{4pt}
\LMDiagramNote{Second takes the Missal and stand to the Epistle side. At the
same time First prepares the folded veil. In this route model, First waits
before crossing A until Second has passed with the book, then continues to
Gospel. Neither carries the chalice or crosses behind the priest.}

\LMAction{1}{Both}{Take the assigned objects.}{First takes the folded chalice
veil at the Epistle side. Second takes the Missal and stand at the Gospel side
with a secure two-handed grip.}
\LMAction{2}{Both}{Change sides and place the objects.}{First places the veil
at its rehearsed Gospel-side point. Second places the Missal and stand straight
on the Epistle side. First gives way at crossing A; keep the marked lanes
and leave the chalice to the priest.}
\EndLMPage
\BeginLMPage{COMMUNION}{FINISH THE EXCHANGE}{%
  \LMTrainerNote{Check}{The veil is at its Gospel point and the Missal stands straight on Epistle. The priest alone handles the chalice.}
  \LMTrainerNote{Route}{Rehearse the two lanes at the actual church. Wait whenever the lane is occupied.}}
\LMPageIntro{Both objects have reached their destinations. Now return empty-handed.}
\LMSceneBookReturnComplete
\LMAction{3}{Both}{Switch back to the normal places.}{After both placements,
genuflect at the centre, pass by the rehearsed lanes, and kneel:
First Epistle side, Second Gospel side.}
\ASGMovementFields{Both assigned objects are placed.}{Both, now empty-handed.}{Meet at the centre; genuflect and pass by the rehearsed lanes.}{Kneel at the normal places on Step 1.}{The priest reassembles and covers the chalice; the servers leave it alone.}
\EndLMPage

% ---------------------------------------------------------------------------
% PAGE 30 — LM-038 through LM-040
% ---------------------------------------------------------------------------
\BeginLMPage{ENDING}{POSTCOMMUNION AND DISMISSAL}{%
  \LMTrainerNote{Prayer over People}{At ferial Masses of Lent and
    Passiontide when appointed, both bow the head at
    \emph{Humiliate capita vestra Deo} and answer Amen at the prayer ending.}
  \LMTrainerNote{Easter}{Card 22 is only for a Mass of the season within the
    Easter octave.}
  \LMTrainerNote{Procession}{If \emph{Benedicamus Domino} announces a
    following liturgical procession, there is no blessing or Last Gospel;
    hand off to that procession's separate sheet.}}
\LMPageIntro{Both kneel at the normal places. Use the familiar responses
unless the Easter-octave dismissal changes the words.}
\LMBranch{Lenten Prayer over the People, when appointed}
{Bow the head at \emph{Humiliate capita vestra Deo}; answer Amen at the
prayer's audible ending, then continue to the dismissal.}
\LMLaneHeadings

\LMDialogueRepeat{12 AGAIN}
  {\emph{Dominus vobiscum.}}
  {A13}

\LMDialogueRepeat{08 AGAIN}
  {\emph{Per omnia saecula saeculorum.}}
  {A9}

\LMDialogueRepeat{12 AGAIN}
  {\emph{Dominus vobiscum.}}
  {A13}

\LMBranch{Outside the Easter octave}{Use the ordinary dismissal below.}
\LMDialogueRepeat{15 AGAIN}
  {\emph{Ite, missa est.}}
  {A15}

\LMBranch{Mass of the season within the Easter octave}
{Use response 22 below instead of the ordinary dismissal.}
\LMDialogueFirst{22}
  {\emph{Ite, missa est, alleluia, alleluia.}}
  {A22}
\LMBranch{If the dismissal is \emph{Benedicamus Domino}}
{Give its appointed response. If it announces a following liturgical
procession, do not continue to the blessing or Last Gospel here; follow the
separate procession route.}
\EndLMPage

% ---------------------------------------------------------------------------
% PAGE 31 — LM-041 and LM-042
% ---------------------------------------------------------------------------
\BeginLMPage{ENDING}{BLESSING AND LAST GOSPEL}{%
  \LMTrainerNote{Blessing}{Both bow at \emph{Benedicat vos omnipotens Deus},
    sign with the blessing, answer Amen, then stand.}
  \LMTrainerNote{Last Gospel}{First, at picture-right on the Epistle side,
    angles toward the priest. Second, at picture-left on the Gospel side,
    faces the altar. Both make the three small crosses at the title.}
  \LMTrainerNote{Omission}{At the third Mass of Christmas and any other
    appointed omission, leave out all three Last-Gospel responses.}}
\LMPageIntro{Kneel and bow the head for the blessing; make the Sign of the
Cross and answer Amen. Then stand. At the Last Gospel's title, make the three
small crosses on forehead, lips, and breast.}
\LMBranch{If the Last Gospel is omitted}
{After the blessing, omit all three Last-Gospel responses and follow the
trained exit route.}
\LMLaneHeadings

\LMDialogueRepeat{08 AGAIN}
  {\emph{Benedicat vos omnipotens Deus: Pater, et Filius, et Spiritus
  Sanctus.}}
  {A9}

\LMDialogueRepeat{12 AGAIN}
  {\emph{Dominus vobiscum.}}
  {A13}

\LMDialogueRepeat{16 AGAIN}
  {\emph{Initium sancti Evangelii secundum Ioannem.}}
  {A16}

\LMDialogueRepeat{15 AGAIN}
  {\emph{Et Verbum caro factum est, et habitavit in nobis: et vidimus
  gloriam eius, gloriam quasi Unigeniti a Patre, plenum gratiae et
  veritatis.}}
  {A15}
\EndLMPage

% ---------------------------------------------------------------------------
% PAGE 32 — LM-042 and LM-043
% ---------------------------------------------------------------------------
\BeginLMPage{ENDING}{FACE LEFT FOR THE LAST GOSPEL}{%
  \LMTrainerNote{Priest}{The priest stands at the Gospel corner and faces left to read. The ordinary route leaves the Missal at Epistle.}
  \LMTrainerNote{Et Verbum}{Both genuflect with the priest at \emph{Et Verbum caro factum est}, then stand again. The picture shows the standing reading state.}
  \LMTrainerNote{Servers}{First turns toward the priest from Epistle; Second remains at Gospel facing the altar.}}
\LMSceneLastGospel
\LMRouteFields{After the blessing, the priest goes to the Gospel corner.}{Both, empty-handed.}{The normal places.}{Stand; First turns toward the priest, while Second faces the altar.}{Both stay at their normal places for the reading and its appointed genuflection.}{The covered chalice remains at the centre; the Missal remains at Epistle.}
\LMDiagramNote{At the final \emph{Deo gratias}, both go to the foot and leave the centre clear. The priest takes his covered chalice and descends.}
\EndLMPage

\BeginLMPage{ENDING}{RECESS TO THE SACRISTY}{%
  \LMTrainerNote{Objects}{The priest carries the covered chalice. If he uses a biretta, First retrieves it and presents it after the final reverence.}
  \LMTrainerNote{Route}{The pictured leftward line is one rehearsed departure direction. Use the real sacristy doorway and turns.}
  \LMTrainerNote{After Mass}{Leonine prayers are outside this Order. If said, use the separately trained local sequence.}}
\LMSceneExit
\LMSceneRecess
\LMRouteFields{The Last Gospel is finished, or its appointed omission applies.}{First, Second, then priest with his covered chalice.}{The stationary formation at the foot.}{Genuflect together on the floor; turn inward and depart single file by the rehearsed sacristy route.}{First leads, Second follows, and the priest follows both; bow together to the sacristy cross.}{After arrival report spills or damage, do only the assigned tidying, and hang up surplice and cassock.}
\EndLMPage

% ---------------------------------------------------------------------------
% TERMINAL APPENDICES — preparation, rehearsal, scope, and sources
% ---------------------------------------------------------------------------
\BeginLMSingleScenePage{APPENDIX}{PREPARE AND MAP YOUR CHURCH}{%
  \LMTrainerNote{First practice}{Walk the complete route empty-handed. Add
    the words. Add cold objects only after the route is calm.}
  \LMTrainerNote{Names}{Use First Acolyte, Second Acolyte, and Both
    consistently throughout rehearsal.}
  \LMTrainerNote{Safety}{A responsible adult controls candles, unfamiliar
    steps, sacred vessels, and every Communion-plate rehearsal.}}
\LMPageIntro{Complete this page with the priest or trainer before practising
with objects. Altar-position scenes are viewed from the nave; close-ups show
hand actions. The real steps, distances, and Communion lane may differ.}
\LMModelKeyLine
\vspace{4pt}
\LMSceneStart
\vspace{2pt}
\LMDiagramNote{\textbf{Opening position:} First stands on the priest's right
(picture-right); Second stands on his left (picture-left). Leave the centre
open.}
\LMAction{1}{Both}{Point to every working place.}{Doorway; foot; normal
places; credence; both Missal corners; Communion lane. Ask before touching
anything unfamiliar.}
\LMAction{2}{Trainer}{Record the local choices.}{Centre stopping point:
\rule{4.2em}{0.5pt}\quad Communion side: \rule{3.8em}{0.5pt}\quad bell
pattern: \rule{3.8em}{0.5pt}.}
\LMBranch{Role summary}
{First begins on the Epistle side and handles the bell, wine, both
post-Communion cruets, towel, Communion plate, Epistle-to-Gospel Missal move,
and post-ablution folded chalice veil. Second begins on the Gospel side and
handles the Offertory water, Lavabo water and bowl, and Gospel-to-Epistle
Missal move.}
\EndLMPage

\BeginLMPage{APPENDIX}{REHEARSAL CHECKLIST}{%
  \LMTrainerNote{Method}{Run each item until both acolytes can complete it
    calmly. Stop and reset after a missed cue rather than rushing onward.}
  \LMTrainerNote{Objects}{Use cold practice objects until handoffs and lanes
    are reliable.}}
\LMPageIntro{These checks gather the practice tests outside the
sacristy-to-sacristy route.}
\LMAction{1}{Both}{Entrance and opening.}{Walk entrance, stop, reverence,
wait, kneel, and answer from \emph{Introibo} through the move to the normal
places. Both finish together and leave the centre open.}
\LMAction{2}{Both}{Readings and Offertory.}{Run the Missal move, Gospel
stations, cruets, Lavabo, bell pickup, and \emph{Orate, fratres} response.
Every object finishes at its named place.}
\LMAction{3}{Both}{Canon.}{With an empty bell, run warning, ascent, both
elevations, and return. No one crosses behind the priest or rings while
walking.}
\LMAction{4}{Both}{Communion.}{Run both branches. First alone serves the
wine, then wine-and-water ablutions and finishes any used plate. Then rehearse
First's veil route and Second's book route together, followed by the switch
back to normal places.}
\LMAction{5}{Both}{Ending.}{Run dismissal, blessing, Last Gospel, final
reverence, and recession. Repeat the no-Last-Gospel branch separately.}
\EndLMPage

\BeginLMPage{APPENDIX}{SCOPE, QUALIFICATIONS, AND SOURCES}{%
  \LMTrainerNote{Use}{This training aid does not replace the Missal, the
    priest's direction, or the rehearsed customs of the actual church.}
  \LMTrainerNote{Flash cards}{Print the paired flash cards on letter paper at
    100 percent, two-sided, flipped on the long edge. Verify one sheet for
    front/back alignment before cutting. Practise with the priest's complete
    cue on the front and Both acolytes' complete response on the back.}
  \LMTrainerNote{Rights}{Received liturgical text and embedded fonts retain
    their external status; see \texttt{THIRD\_PARTY.md}.}}
\LMPageIntro{This guide covers ordinary two-server Low Mass under the 1962
\emph{Missale Romanum}. Its three-level sanctuary, named First/Second
division, and practical routes are documented teaching models, not universal
rubrics. It teaches the Italianate ecclesiastical-Latin pronunciation recorded
in the series audit; its sound lines and meanings are project-written learning
aids, not text to say at Mass. Special rites, Requiem Mass, Holy Week, exposed-
Sacrament service, adjacent devotions, and the Leonine prayers require separate
guidance. Present permission to use the 1962 books and authorization to serve
belong to the competent authorities; this guide grants neither.}
\LMBranch{Custom and rights}
{Where the Missal does not fix a practical route, this guide follows the named
manual model or a choice rehearsed for the actual church. Received liturgical
text, source quotations, fonts, and other third-party material retain their
external rights status and are not relicensed by the project.}
\LMBranch{Principal references}
{\emph{Missale Romanum} (1962), especially \emph{Ritus servandus} II--XII
and General Rubrics 424--425 and 503--510; Archconfraternity of Saint
Stephen, \emph{Serving Low Mass with 2 Servers (Ac1 \& Ac2)}; Charles
J. Carmody, \emph{Learning to Serve} (1961); Laurence J. O'Connell,
\emph{The Book of Ceremonies} (1943), pp.~170--174; and the 1961
\emph{Liber Usualis} pronunciation rules.}
\par\vspace{10pt}
\begingroup\RaggedRight
\ifChildBookletGuide
  % The child lane is full width, so its five required provenance records
  % need a compact terminal display to leave the common 7pt rights colophon
  % wholly inside the physical page.
  \fontsize{8.4}{9.7}\selectfont
\fi
\LMRenderGenerationMetadata\par
\endgroup
\TriptychRightsNotice
\EndLMPage
